\section{Questão 1}

Considere um processo randômico em tempo discreto $X[n]$ definido como
%
$$
    X[n] = A\cos(\omega_0 n + \Phi)
$$
%
onde $u(t)$ é a função degrau unitário,
a frequência digital $\omega_0 = 3\frac{\pi}{16}$ (radianos/amostra),
$A>0$, $A\in\mathbb{R}$ é uma constante e 
$\Phi \sim \text{Uniform}[0, 2\pi]$, isto é,
$\Phi$ tem distribuição de probabilidade uniforme e o espaço amostral de $\Phi$ é $(0, 2\pi)$.

\begin{itemize}
    \item [a)] Encontre a função da expectância $\mu[n]$.
    \item [b)] Encontre a função de correlação de $R[m,n]$.
    \item [c)] Este processo é randômico ou WSS? Justifique.
\end{itemize}

\begin{respostas}

\begin{itemize}
    \item [a)] 
    Como a cossenoide tem média zero em um período, espera-se que $\mu[n] = 0$.
    De fato, tem-se
    %
    \begin{align}
        \mu[n] &= \E{X[n]} \\
        &= \E{A\cos(\omega_0 n + \Phi)} \\
        &= A\E{\cos(\omega_0 n + \Phi)} \\
        &= A\int_{0}^{2\pi} \cos(\omega_0 n + \phi) f_\Phi(\phi) d\phi \\
        &= A\int_{0}^{2\pi} \cos(\omega_0 n + \phi) \frac{1}{2\pi} d\phi \\
        &= 0.
    \end{align}
    \item [b)] Tem-se
    %
    \begin{align}
        R[m,n] &= \E{X[m]X[n]} \\
        &= \E{A\cos(\omega_0 m + \Phi)A\cos(\omega_0 n + \Phi)} \\
        &= A^2\E{\cos(\omega_0 m + \Phi)\cos(\omega_0 n + \Phi)} \\
        &= A^2\int_{0}^{2\pi} \cos(\omega_0 m + \phi)\cos(\omega_0 n + \phi) f_\Phi(\phi) d\phi \\
        &= A^2\int_{0}^{2\pi} \cos(\omega_0 m + \phi)\cos(\omega_0 n + \phi) \frac{1}{2\pi} d\phi \\
        &= \frac{A^2}{2\pi} \int_{0}^{2\pi} \frac{1}{2} [\cos(\omega_0 (m-n)) + \cos(\omega_0 (m+n) + 2\phi)] d\phi \\
        &= \frac{A^2}{4\pi} \left[
                \int_{0}^{2\pi} \cos(\omega_0 (m-n)) d\phi + 
                \int_{0}^{2\pi} \cos(\omega_0 (m+n) + 2\phi) d\phi 
            \right] \\
        &= \frac{A^2}{4\pi} \left[
                2\pi \cos(\omega_0 (m-n)) + 0
            \right] \\
        &= \frac{A^2}{2} \cos(\omega_0 (m-n)).
    \end{align}

    \item [c)] $R[m,n]$ depende apenas da diferença $m-n$, tornando o processo WSS.
\end{itemize}

\end{respostas}